% Partie : set-theory
% Chapitre : spine
% Section : rank

\documentclass[../../../include/open-logic-section]{subfiles}

\begin{document}

\olfileid[fr]{sth}{spine}{rank}
\olsection{Rang}

Maintenant que nous avons défini les étapes comme les $V_\alpha$ et
que nous savons que tout ensemble est un sous-ensemble d'une certaine
étape, nous pouvons définir le \emph{rang} d'un ensemble.
Intuitivement, le rang de $A$ est le premier moment où $A$ se forme.
Plus précisément :

\begin{defn}\ollabel{defnsetrank}
Pour chaque ensemble $A$, $\setrank{A}$ est le plus petit ordinal
$\alpha$ tel que $A \subseteq V_\alpha$.
\end{defn}
\begin{prop}\ollabel{ranksexist}
$\setrank{A}$ existe pour tout ensemble $A$.
\end{prop}
\begin{proof}
La démonstration est laissée en exercice.
\end{proof}
\begin{prob}
Démontrez la \olref[sth][spine][rank]{ranksexist}.
\end{prob}
\noindent
Le bon ordre des rangs permet d'établir quelques résultats importants :
\begin{prop}\ollabel{valphalowerrank}
Pour tout ordinal $\alpha$,
$V_\alpha = \Setabs{x}{\setrank{x} \in \alpha}$.
\end{prop}

\begin{proof}
Si $\setrank{x} \in \alpha$, alors
$x \subseteq V_{\setrank{x}} \in V_\alpha$ ; donc
$x \in V_\alpha$ puisque $V_\alpha$ est potent (en appliquant
plusieurs fois le \olref[Valphabasic]{Valphabasicprops}).
Réciproquement, si $x \in V_\alpha$, alors
$x \subseteq V_\alpha$, donc $\setrank{x} \leq \alpha$. Une
récurrence transfinie simple montre en fait que
$\setrank{x} \in \alpha$.\footnote{Note de l'édition française :
l'original conclut ici que $x \notin V_\alpha$, ce qui contredit
l'hypothèse $x \in V_\alpha$. La récurrence doit établir
$\setrank{x} \in \alpha$ ; ses trois cas sont laissés à l'exercice
suivant.}
\end{proof}
\begin{prob}
Complétez la récurrence transfinie simple mentionnée dans la
\olref[sth][spine][rank]{valphalowerrank}.
\end{prob}

\begin{prop}\ollabel{rankmemberslower}
Si $B \in A$, alors $\setrank{B} \in \setrank{A}$.
\end{prop}

\begin{proof}
$A \subseteq V_{\setrank{A}} =
\Setabs{x}{\setrank{x} \in \setrank{A}}$ par la
\olref{valphalowerrank}.
\end{proof}
\noindent
Ce fait permet d'établir un résultat grâce auquel nous pourrons
démontrer des propriétés de \emph{tous les ensembles} par une forme
de récurrence :

\begin{thm}[Schéma de récurrence sur $\in$]
Pour toute formule $\phi$ :
\[
	\forall A((\forall x \in A)\phi(x) \lif \phi(A)) \lif \forall A \phi(A).
\]
\end{thm}

\begin{proof}
Nous démontrons la contraposée. Supposons donc
$\lnot \forall A \phi(A)$. Par récurrence transfinie
(\olref[ordinals][basic]{ordinductionschema}), il existe un ensemble
qui ne vérifie pas $\phi$ et dont le rang est le plus petit possible :
un ensemble $A$ tel que $\lnot \phi(A)$ et
$\forall x(\setrank{x} \in \setrank{A} \lif \phi(x))$. Or, si
$x \in A$, alors $\setrank{x} \in \setrank{A}$ par la
\olref{rankmemberslower}, donc $\phi(x)$. Ainsi,
$(\forall x \in A)\phi(x) \land \lnot \phi(A)$.
\end{proof}\noindent On peut présenter ce puissant résultat de manière
informelle. Disons que $\phi$ est \emph{héréditaire} si, dès que tous
les !!{element}s d'un ensemble vérifient $\phi$, l'ensemble lui-même
la vérifie. La récurrence sur $\in$ dit alors : si $\phi$ est
héréditaire, tout ensemble vérifie $\phi$.

Pour terminer, pour l'instant, notre étude des rangs, démontrons
quelques affirmations annoncées plus haut.

\begin{prop}\ollabel{ranksupstrict}
$\setrank{A} = \supstrict_{x \in A}\setrank{x}$.
\end{prop}

\begin{proof}
Posons $\alpha = \supstrict_{x \in A}\setrank{x}$. Par la
\olref{rankmemberslower}, $\alpha \leq \setrank{A}$. Mais, si
$x \in A$, alors $\setrank{x} \in \alpha$ ; donc
$x \in V_\alpha$ par la \olref{valphalowerrank}, et par suite
$A \subseteq V_\alpha$, c'est-à-dire
$\setrank{A} \leq \alpha$. Ainsi, $\setrank{A} = \alpha$.
\end{proof}

\begin{cor}\ollabel{ordsetrankalpha}
Pour tout ordinal $\alpha$, $\setrank{\alpha} = \alpha$.
\end{cor}

\begin{proof}
Comme hypothèse de récurrence transfinie, supposons que
$\setrank{\beta} = \beta$ pour tout $\beta \in \alpha$. Alors
$\setrank{\alpha} = \supstrict_{\beta \in \alpha}\setrank{\beta}
= \supstrict_{\beta \in \alpha}\beta = \alpha$ par la
\olref{ranksupstrict}.
% Notons d'abord que \setrank{\alpha} \neq \beta pour tout
% \beta \in \alpha. Sinon, \beta = \setrank{\beta} \in
% \setrank{\alpha} = \beta par \olref{rankmemberslower}, soit
% \beta \in \beta, contradiction.
%
% Notons ensuite que \alpha \subseteq V_\alpha. En effet,
% \alpha = \Setabs{\beta}{\beta \in \alpha} par
% \olref[sfr][ordinals][basic]{ordissetofsmallerord}. Si
% \beta \in \alpha, alors \beta \subseteq V_{\setrank{\beta}} =
% V_\beta \in V_\alpha, donc \beta \in V_\alpha par
% \olref[sfr][spine][Valphabasic]{Valphabasicprops} deux fois.
\end{proof}

Enfin, voici une courte démonstration du résultat promis à la fin de la
\olref[foundation]{sec} : $\ZFminus$ démontre l'implication
$\emph{Régularité} \Rightarrow \emph{Fondation}$. (Notons que la
notion de «~rang~» et la \olref{rankmemberslower} peuvent servir ici :
comme nous l'avons indiqué au début de cette section, ils sont
disponibles dans $\ZFminus + \text{Régularité}$.)

\begin{prop}[dans $\ZFminus + \text{Régularité}$]\ollabel{zfminusregularityfoundation}
La fondation est vraie.
\end{prop}

\begin{proof}
Fixons $A \neq \emptyset$ et choisissons un
$B \in A$ de rang aussi petit que possible. Si $c \in B$, alors
$\setrank{c} \in \setrank{B}$ par la \olref{rankmemberslower} ; donc
$c \notin A$ par le choix de $B$.
\end{proof}

\end{document}
